\subsection{Continuous functions vanishing at infinity on a locally compact space}

In this section, X is a locally compact space. We denote by \(\mathcal{C}_{0}(X)\) the commutative Banach algebra of continuous complex-valued functions vanishing at infinity on X, equipped with the norm \(\|f\|=\sup_{x\in X}|f(x)|\) (example 3 of I, p.~17).

PROPOSITION 1. --- For every closed subset \(\Phi\) of X, let \(I_{\Phi}\) be the set of \(f \in \mathcal{C}_0(X)\) vanishing on \(\Phi\) . Then \(\Phi \mapsto I_{\Phi}\) is a bijection from the set of closed subsets of X onto the set of closed ideals of \(\mathcal{C}_0(X)\) .

The set \(\mathrm{I}_{\Phi}\) is a closed ideal of \(\mathcal{C}_0(\mathrm{X})\) .

Let \(\Phi \neq \Phi'\) be closed subsets of X. If necessary, by exchanging \(\Phi\) and \(\Phi'\) one may assume that there exists \(x \in \Phi'\) such that \(x \notin \Phi\) ; there then exists a function \(f \in \mathcal{C}_0(X)\) vanishing on \(\Phi\) and nonzero at \(x\) (TG, IX, p.~43, prop. 1). We have \(f \in I_{\Phi}\) and \(f \notin I_{\Phi'}\) , so that the map \(\Phi \mapsto I_{\Phi}\) is injective.

Let I be a closed ideal of \(\mathcal{C}_{0}(\mathrm{X})\) . Let \(\Phi\) be the set of \(x \in X\) such that \(f(x) = 0\) for all \(f \in I\) ; it is a closed subset of X, and we have \(I \subset I_{\Phi}\) . Let us prove that \(I_{\Phi} \subset I\) , which will imply that \(I = I_{\Phi}\) and will finish the proof of the proposition.

Let \(f \in \mathrm{I}_{\Phi}\) . For every real number \(\varepsilon > 0\) , let \(\mathrm{C}_{\varepsilon}\) be the set of \(x \in \mathrm{X}\) such that \(|f(x)| \geqslant \varepsilon\) . Since \(f\) tends to 0 at infinity, the set \(\mathrm{C}_{\varepsilon}\) is compact. Let \(x \in \mathrm{C}_{\varepsilon}\) ; since \(f(x) \neq 0\) and \(f \in \mathrm{I}_{\Phi}\) , one has \(x \notin \Phi\) ; by definition of \(\Phi\) , there then exists a function \(\varphi_x \in \mathrm{I}\) such that \(|\varphi_x(x)| > 1\) , hence such that \(|\varphi_x(y)| > 1\) for all \(y\) belonging to a neighborhood \(\mathrm{V}_x\) of \(x\) . The open sets \(\mathrm{V}_x \cap \mathrm{C}_{\varepsilon}\) cover \(\mathrm{C}_{\varepsilon}\) . Since the set \(\mathrm{C}_{\varepsilon}\) is compact, there exists a finite subset \(\mathrm{T}_{\varepsilon} \subset \mathrm{X}\) such that

\[
\mathrm{C} _ {\varepsilon} \subset \bigcup_ {x \in \mathrm{T} _ {\varepsilon}} \mathrm{V} _ {x}.
\]

Then the element

\[
g _ {\varepsilon} = \frac {1}{\varepsilon} \sum_ {x \in \mathrm{T} _ {\varepsilon}} \varphi_ {x} \overline {{\varphi_ {x}}} \geqslant 0
\]

of \(\mathcal{C}_0(\mathrm{X})\) belongs to I, and we have \(g_{\varepsilon} \geqslant \varepsilon^{-1}\) on \(\mathrm{C}_{\varepsilon}\) . The function

\[
f _ {\varepsilon} = \frac {f g _ {\varepsilon}}{1 + g _ {\varepsilon}}
\]

belongs to I. For \(x \notin \mathrm{C}_{\varepsilon}\) , one has

\[
\left| f (x) - f _ {\varepsilon} (x) \right| \leqslant 2 \varepsilon ,
\]

and for \(x\in \mathrm{C}_{\varepsilon}\) , one has

\[
| f (x) - f _ {\varepsilon} (x) | = \frac {| f (x) |}{1 + g _ {\varepsilon} (x)} \leqslant \varepsilon | f (x) |.
\]

Thus \(f_{\varepsilon}\) converges uniformly to f on X as \(\varepsilon\) tends to 0. We therefore have \(f \in \bar{I}\) , whence \(f \in I\) since I is closed.

COROLLARY 1. --- For every \(x \in X\) , let \(I_x\) be the set of \(f \in \mathcal{C}_0(X)\) vanishing at \(x\) . Then \(x \mapsto I_x\) is a bijection from \(X\) onto the set of maximal closed ideals of \(\mathcal{C}_0(X)\) . These ideals are regular.

This follows immediately from prop. 1.

Let X' denote the Alexandroff compactification of X, that is, the compact space obtained from X by adjoining a point at infinity \(\omega_{X}\) (TG, I, pp. 67–68). The algebra \(\mathcal{C}_{0}(X)\) is identified with the Banach algebra of continuous complex-valued functions on X' vanishing at \(\omega_{X}\) .

For every \(x \in X'\) , we denote by \(ev_{x}\) the character of \(\mathcal{C}_{0}(X)\) defined by \(\mathrm{ev}_{x}(f) = f(x)\) for all \(f \in \mathcal{C}_{0}(X)\) .

COROLLARY 2. --- The map \(x \mapsto \mathrm{ev}_x\) is a homeomorphism from \(\mathbf{X}'\) onto \(\mathsf{X}'(\mathcal{C}_0(\mathbf{X}))\) and its restriction to \(\mathbf{X}\) is a homeomorphism from \(\mathbf{X}\) onto \(\mathsf{X}(\mathcal{C}_0(\mathbf{X}))\) . Moreover, the weak topology and the Jacobson topology coincide on \(\mathsf{X}(\mathcal{C}_0(\mathbf{X}))\) .

The map ev: \(x \mapsto ev_{x}\) of \(X'\) into \(\mathsf{X}'(\mathcal{C}_{0}(\mathbf{X}))\) is injective. It is surjective by Cor. 1 and Th. 2 of I, p.~30. It is continuous, because for every function \(f \in \mathcal{C}_{0}(\mathbf{X})\) and every open set U of R, we have

\[
\mathrm{ev} ^ {- 1} (\{\chi \in \mathsf {X} ^ {\prime} (\mathcal {C} _ {0} (\mathsf {X})) \mid \chi (f) \in \mathrm{U} \}) = f ^ {- 1} (\mathrm{U})
\]

which is open in X. The map ev is therefore a homeomorphism since \(X'\) is compact. The restriction of ev to X is then a homeomorphism onto \(\mathsf{X}(\mathcal{C}_{0}(\mathsf{X}))\) .

If F is a weakly closed subset of \(\mathsf{X}(\mathcal{C}_{0}(\mathbf{X}))\) it corresponds, via the homeomorphism ev, to a closed subset \(\Phi\) of X; precisely, by Prop. 1, we have \(\mathrm{F} = \{\chi \in \mathsf{X}(\mathcal{C}_{0}(\mathbf{X})) \mid \mathrm{I}_{\Phi} \subset \mathrm{Ker} \, \chi\}\) which is closed for the Jacobson topology.

COROLLARY 3. --- Suppose X is compact. Then the map \(x \mapsto \mathrm{ev}_x\) is a homeomorphism from X onto \(\mathsf{X}(\mathcal{C}(\mathbb{X}))\) . The weak topology and the Jacobson topology coincide on \(\mathsf{X}(\mathcal{C}(\mathbb{X}))\) .
% semantic-fragments: legacy-input-seam
\relax{}
